\section{Methods}
\paragraph{Member-derived KC representations.}
Let $q_i$ be a question text assigned to KC $c$, and let $e(q_i)$ be its embedding. A member-derived representation for $c$ is an aggregation over $\{e(q_i): q_i \in c\}$. We instantiate this with mean, median, and medoid prototypes where artifacts permit. Mean prototypes are the primary condition because they correspond directly to the common class-prototype assumption: a category can be summarized by the central tendency of its members.

\paragraph{Category-validity analyses.}
We first test whether prototypes recover the human KC labels. For each question, we retrieve the nearest KC prototype and compute Recall@1, Recall@5, MRR, NDCG, and a shuffled-prototype control. This is an optimistic diagnostic when the queried question contributes to its own KC prototype, so we interpret it as a category-signal check rather than a deployable classifier. We also measure coherence by comparing intra-KC and inter-KC cosine similarity and recording cluster diagnostics. Finally, we test whether prototype geometry aligns with coarse label information extracted heuristically from KC names. Because explicit curriculum metadata are incomplete, these label-alignment results are treated as weak supervision rather than gold taxonomies.

\paragraph{Functional analyses.}
Category recovery does not establish functional validity. We therefore evaluate three non-DKT downstream functionalities before the DKT initialization test. First, we compare prototype similarity graphs with KT-relevant graphs: transition relations, co-occurrence, difficulty similarity, and available DKT-derived KC relation artifacts. Second, we predict item difficulty from representations, treating difficulty as a psychometric functionality that a useful educational label representation might encode. Third, we predict KC transitions, treating next-KC structure as a sequential functionality of the label space. These diagnostics are not substitutes for DKT, but they test whether the representations encode behavioral structure beyond label membership.

\paragraph{Chapter 5 downstream test.}
The main downstream evidence is Chapter 5. DKT consumes KC-response inputs, so Chapter 5 constructs response-specific embeddings through nine pipelines. The pipelines vary whether question text is aggregated into KCs before adding response state, whether response state is added before aggregation, and whether KC-level or question-level text windows are used. We consolidate four Chapter 5 result families: full-data DKT initialization, reduced-data DKT initialization, fine-tuned BERT initialization, and post hoc alignment between semantic KC graphs and DKT-derived KC graphs. This paper treats those experiments as the decisive test of whether member-derived category representations become useful functional priors.
